% --- Archon physics formalization source begin ---
% archon:physics
% archon:covers PhyXMiniProblems/problem_phyx_mini_0219.lean
% archon:source-report reports/phyx_mini/problem_phyx_mini_0219.source.json

\chapter{Physics problem phyx\_mini\_0219}
\label{ch:PhyXMiniProblems_problem_phyx_mini_0219}

\paragraph{Problem source.}
\#\# Physical scenario

Two loudspeakers are placed \textbackslash{}( 3.00\textbackslash{},\textbackslash{}mathrm\{m\} \textbackslash{}) apart, as shown in figure. They emit \textbackslash{}( 474\textbackslash{},\textbackslash{}mathrm\{Hz\} \textbackslash{}) sounds, in phase. A microphone is placed \textbackslash{}( 3.20\textbackslash{},\textbackslash{}mathrm\{m\} \textbackslash{}) distant from a point midway between the two speakers, where an intensity maximum is recorded.

\#\# Question

How far must the microphone be moved to the right to find the first intensity minimum?

\#\# Answer choices

- A: A: 0.729m

- B: B: 0.629m

- C: C: 0.529m

- D: D: 0.429m

\#\# Figure caption (auxiliary; use the image as primary evidence)

The image illustrates a geometric arrangement involving two speakers and a microphone. There's a light blue background. At the top, two speakers are shown, each with a wire connecting to its top. The speakers are positioned horizontally in line with each other, separated by a labeled distance of "3.00 m."

Below the speakers, there is a microphone, indicated at the point of the triangle formed by dashed lines extending from each speaker, labeled as \textbackslash{}(d\_1\textbackslash{}) and \textbackslash{}(d\_2\textbackslash{}). The distance between each speaker and the microphone is labeled as "3.20 m." 

The dashed lines form a triangular shape with the microphone at the vertex, illustrating the spatial relationship between the speakers and the microphone.

\#\# Dataset metadata

Resonance and Harmonics; Physical Model Grounding Reasoning, Spatial Relation Reasoning

\paragraph{Recorded answer/context.}
D: D: 0.429m

\paragraph{Figure/image path.}
/root/proposal\_for\_physic/science-mango/phyx\_mini\_run/phyx\_data/test\_image/219.png

\paragraph{Formalization target.}
create a compiling Lean file with sorry bodies at `PhyXMiniProblems/problem\_phyx\_mini\_0219.lean`. The Lean declarations must preserve the physical quantities, dimensions or dimensional roles, figure labels, governing-law hypotheses, and final relation expressed by this problem.
Use Mathlib/Physlib names found through LeanExplore where available. If a physics API is missing, introduce faithful local abstractions rather than scalar placeholder aliases.

\begin{theorem}[Physics formalization target]
\label{thm:physics:phyx_mini_0219:target}
\lean{PhyXMiniProblems.ProblemPhyXMini0219.problem_phyx_mini_0219}
\uses{def:physics:phyx-mini-0219:phyxminiproblems-problemphyxmini0219-lengthquantity, def:physics:phyx-mini-0219:phyxminiproblems-problemphyxmini0219-twospeakerinterferencesetup, def:physics:phyx-mini-0219:phyxminiproblems-problemphyxmini0219-matchescoherentinphasesourcescenario, def:physics:phyx-mini-0219:phyxminiproblems-problemphyxmini0219-matchesproblemreadouts, def:physics:phyx-mini-0219:phyxminiproblems-problemphyxmini0219-matchessuppliedfigure, def:physics:phyx-mini-0219:phyxminiproblems-problemphyxmini0219-usesstandardroomairsoundspeed, def:physics:phyx-mini-0219:phyxminiproblems-problemphyxmini0219-hasphysicalparameters, def:physics:phyx-mini-0219:phyxminiproblems-problemphyxmini0219-satisfiesacousticwavelaw, def:physics:phyx-mini-0219:phyxminiproblems-problemphyxmini0219-satisfiesrightmoveandpathgeometry, def:physics:phyx-mini-0219:phyxminiproblems-problemphyxmini0219-satisfiesinphasedestructiveinterferencelaw, def:physics:phyx-mini-0219:phyxminiproblems-problemphyxmini0219-isfirstpositiverightmovetominimum, def:physics:phyx-mini-0219:phyxminiproblems-problemphyxmini0219-recordeddatasetanswer, def:physics:phyx-mini-0219:phyxminiproblems-problemphyxmini0219-roundstodisplayedmove, def:physics:phyx-mini-0219:phyxminiproblems-problemphyxmini0219-isuniquematchinganswer}
The assigned autoformalize agent should translate this physics problem into Lean declarations in the covered file, with theorem and lemma proof bodies written as `by sorry`.
\end{theorem}
\begin{proof}
This is an autoformalization task, not a proof task. Produce statements that can later be proved by the physics prover without weakening the physical model.
\end{proof}

% --- Archon declaration topology begin ---
\section{Declaration topology}
\begin{definition}[Length Quantity]
  \label{def:physics:phyx-mini-0219:phyxminiproblems-problemphyxmini0219-lengthquantity}
  \lean{PhyXMiniProblems.ProblemPhyXMini0219.LengthQuantity}
  A nonnegative physical length independent of a chosen readout unit
\end{definition}
\begin{proof}
  This is the declared model definition of Length Quantity.
\end{proof}
\begin{definition}[Signed Length Quantity]
  \label{def:physics:phyx-mini-0219:phyxminiproblems-problemphyxmini0219-signedlengthquantity}
  \lean{PhyXMiniProblems.ProblemPhyXMini0219.SignedLengthQuantity}
  A signed physical length used for Cartesian figure coordinates
\end{definition}
\begin{proof}
  This is the declared model definition of Signed Length Quantity.
\end{proof}
\begin{definition}[Frequency Quantity]
  \label{def:physics:phyx-mini-0219:phyxminiproblems-problemphyxmini0219-frequencyquantity}
  \lean{PhyXMiniProblems.ProblemPhyXMini0219.FrequencyQuantity}
  A nonnegative physical frequency carrying inverse-time dimension
\end{definition}
\begin{proof}
  This is the declared model definition of Frequency Quantity.
\end{proof}
\begin{definition}[Speed Quantity]
  \label{def:physics:phyx-mini-0219:phyxminiproblems-problemphyxmini0219-speedquantity}
  \lean{PhyXMiniProblems.ProblemPhyXMini0219.SpeedQuantity}
  A nonnegative physical sound speed
\end{definition}
\begin{proof}
  This is the declared model definition of Speed Quantity.
\end{proof}
\begin{definition}[Acoustic Intensity Quantity]
  \label{def:physics:phyx-mini-0219:phyxminiproblems-problemphyxmini0219-acousticintensityquantity}
  \lean{PhyXMiniProblems.ProblemPhyXMini0219.AcousticIntensityQuantity}
  Acoustic intensity has dimension mass per time cubed, equivalently power per area. It is kept as a physical quantity even though only extrema, rather than numerical intensity readouts, are used by this problem
\end{definition}
\begin{proof}
  This is the declared model definition of Acoustic Intensity Quantity.
\end{proof}
\begin{definition}[length Readout]
  \label{def:physics:phyx-mini-0219:phyxminiproblems-problemphyxmini0219-lengthreadout}
  \lean{PhyXMiniProblems.ProblemPhyXMini0219.lengthReadout}
  \uses{def:physics:phyx-mini-0219:phyxminiproblems-problemphyxmini0219-lengthquantity}
  Read a physical length in a selected length unit
\end{definition}
\begin{proof}
  This is the declared model definition of length Readout.
\end{proof}
\begin{definition}[signed Length Readout]
  \label{def:physics:phyx-mini-0219:phyxminiproblems-problemphyxmini0219-signedlengthreadout}
  \lean{PhyXMiniProblems.ProblemPhyXMini0219.signedLengthReadout}
  \uses{def:physics:phyx-mini-0219:phyxminiproblems-problemphyxmini0219-signedlengthquantity}
  Read a signed Cartesian coordinate in a selected length unit
\end{definition}
\begin{proof}
  This is the declared model definition of signed Length Readout.
\end{proof}
\begin{definition}[frequency Readout]
  \label{def:physics:phyx-mini-0219:phyxminiproblems-problemphyxmini0219-frequencyreadout}
  \lean{PhyXMiniProblems.ProblemPhyXMini0219.frequencyReadout}
  \uses{def:physics:phyx-mini-0219:phyxminiproblems-problemphyxmini0219-frequencyquantity}
  Read frequency in inverse units of a selected time unit
\end{definition}
\begin{proof}
  This is the declared model definition of frequency Readout.
\end{proof}
\begin{definition}[speed Readout]
  \label{def:physics:phyx-mini-0219:phyxminiproblems-problemphyxmini0219-speedreadout}
  \lean{PhyXMiniProblems.ProblemPhyXMini0219.speedReadout}
  \uses{def:physics:phyx-mini-0219:phyxminiproblems-problemphyxmini0219-speedquantity}
  Read speed in a selected length unit per selected time unit
\end{definition}
\begin{proof}
  This is the declared model definition of speed Readout.
\end{proof}
\begin{definition}[length In Meters]
  \label{def:physics:phyx-mini-0219:phyxminiproblems-problemphyxmini0219-lengthinmeters}
  \lean{PhyXMiniProblems.ProblemPhyXMini0219.lengthInMeters}
  \uses{def:physics:phyx-mini-0219:phyxminiproblems-problemphyxmini0219-lengthquantity, def:physics:phyx-mini-0219:phyxminiproblems-problemphyxmini0219-lengthreadout}
  Meter readout of a nonnegative physical length
\end{definition}
\begin{proof}
  This is the declared model definition of length In Meters.
\end{proof}
\begin{definition}[signed Length In Meters]
  \label{def:physics:phyx-mini-0219:phyxminiproblems-problemphyxmini0219-signedlengthinmeters}
  \lean{PhyXMiniProblems.ProblemPhyXMini0219.signedLengthInMeters}
  \uses{def:physics:phyx-mini-0219:phyxminiproblems-problemphyxmini0219-signedlengthquantity, def:physics:phyx-mini-0219:phyxminiproblems-problemphyxmini0219-signedlengthreadout}
  Meter readout of a signed figure coordinate
\end{definition}
\begin{proof}
  This is the declared model definition of signed Length In Meters.
\end{proof}
\begin{definition}[frequency In Hertz]
  \label{def:physics:phyx-mini-0219:phyxminiproblems-problemphyxmini0219-frequencyinhertz}
  \lean{PhyXMiniProblems.ProblemPhyXMini0219.frequencyInHertz}
  \uses{def:physics:phyx-mini-0219:phyxminiproblems-problemphyxmini0219-frequencyquantity, def:physics:phyx-mini-0219:phyxminiproblems-problemphyxmini0219-frequencyreadout}
  Hertz readout of a physical frequency
\end{definition}
\begin{proof}
  This is the declared model definition of frequency In Hertz.
\end{proof}
\begin{definition}[speed In Meters Per Second]
  \label{def:physics:phyx-mini-0219:phyxminiproblems-problemphyxmini0219-speedinmeterspersecond}
  \lean{PhyXMiniProblems.ProblemPhyXMini0219.speedInMetersPerSecond}
  \uses{def:physics:phyx-mini-0219:phyxminiproblems-problemphyxmini0219-speedquantity, def:physics:phyx-mini-0219:phyxminiproblems-problemphyxmini0219-speedreadout}
  Meter-per-second readout of a physical speed
\end{definition}
\begin{proof}
  This is the declared model definition of speed In Meters Per Second.
\end{proof}
\begin{definition}[Speaker Label]
  \label{def:physics:phyx-mini-0219:phyxminiproblems-problemphyxmini0219-speakerlabel}
  \lean{PhyXMiniProblems.ProblemPhyXMini0219.SpeakerLabel}
  The left and right loudspeakers visible in the primary image
\end{definition}
\begin{proof}
  This is the declared model definition of Speaker Label.
\end{proof}
\begin{definition}[Path Label]
  \label{def:physics:phyx-mini-0219:phyxminiproblems-problemphyxmini0219-pathlabel}
  \lean{PhyXMiniProblems.ProblemPhyXMini0219.PathLabel}
  The two dashed source-to-microphone paths labelled in the image
\end{definition}
\begin{proof}
  This is the declared model definition of Path Label.
\end{proof}
\begin{definition}[speaker For Path]
  \label{def:physics:phyx-mini-0219:phyxminiproblems-problemphyxmini0219-speakerforpath}
  \lean{PhyXMiniProblems.ProblemPhyXMini0219.speakerForPath}
  \uses{def:physics:phyx-mini-0219:phyxminiproblems-problemphyxmini0219-speakerlabel, def:physics:phyx-mini-0219:phyxminiproblems-problemphyxmini0219-pathlabel}
  The image labels d₁ from the right speaker and d₂ from the left
\end{definition}
\begin{proof}
  This is the declared model definition of speaker For Path.
\end{proof}
\begin{definition}[Acoustic Medium]
  \label{def:physics:phyx-mini-0219:phyxminiproblems-problemphyxmini0219-acousticmedium}
  \lean{PhyXMiniProblems.ProblemPhyXMini0219.AcousticMedium}
  The propagation medium implicit in the classroom sound problem
\end{definition}
\begin{proof}
  This is the declared model definition of Acoustic Medium.
\end{proof}
\begin{definition}[Acoustic Source Kind]
  \label{def:physics:phyx-mini-0219:phyxminiproblems-problemphyxmini0219-acousticsourcekind}
  \lean{PhyXMiniProblems.ProblemPhyXMini0219.AcousticSourceKind}
  The kind of coherent acoustic source drawn in the figure
\end{definition}
\begin{proof}
  This is the declared model definition of Acoustic Source Kind.
\end{proof}
\begin{definition}[Interference Approximation]
  \label{def:physics:phyx-mini-0219:phyxminiproblems-problemphyxmini0219-interferenceapproximation}
  \lean{PhyXMiniProblems.ProblemPhyXMini0219.InterferenceApproximation}
  The multiple-choice calculation uses the standard path-difference-only model of two-source interference; distance-dependent amplitude variation is not used to shift the locations of the extrema
\end{definition}
\begin{proof}
  This is the declared model definition of Interference Approximation.
\end{proof}
\begin{definition}[Plane Point]
  \label{def:physics:phyx-mini-0219:phyxminiproblems-problemphyxmini0219-planepoint}
  \lean{PhyXMiniProblems.ProblemPhyXMini0219.PlanePoint}
  \uses{def:physics:phyx-mini-0219:phyxminiproblems-problemphyxmini0219-signedlengthquantity}
  A Cartesian point whose coordinates are dimensionful signed lengths
\end{definition}
\begin{proof}
  This is the declared model definition of Plane Point.
\end{proof}
\begin{definition}[Two Speaker Interference Setup]
  \label{def:physics:phyx-mini-0219:phyxminiproblems-problemphyxmini0219-twospeakerinterferencesetup}
  \lean{PhyXMiniProblems.ProblemPhyXMini0219.TwoSpeakerInterferenceSetup}
  \uses{def:physics:phyx-mini-0219:phyxminiproblems-problemphyxmini0219-lengthquantity, def:physics:phyx-mini-0219:phyxminiproblems-problemphyxmini0219-frequencyquantity, def:physics:phyx-mini-0219:phyxminiproblems-problemphyxmini0219-speedquantity, def:physics:phyx-mini-0219:phyxminiproblems-problemphyxmini0219-acousticintensityquantity, def:physics:phyx-mini-0219:phyxminiproblems-problemphyxmini0219-speakerlabel, def:physics:phyx-mini-0219:phyxminiproblems-problemphyxmini0219-pathlabel, def:physics:phyx-mini-0219:phyxminiproblems-problemphyxmini0219-acousticmedium, def:physics:phyx-mini-0219:phyxminiproblems-problemphyxmini0219-acousticsourcekind, def:physics:phyx-mini-0219:phyxminiproblems-problemphyxmini0219-interferenceapproximation, def:physics:phyx-mini-0219:phyxminiproblems-problemphyxmini0219-planepoint}
  Independent physical data and observables for the pictured experiment. The path lengths d₁ and d₂, the wavelength, and the locations of intensity minima are related only by the governing-law structures below. In particular, no field assigns the requested movement a numerical value
\end{definition}
\begin{proof}
  This is the declared model definition of Two Speaker Interference Setup.
\end{proof}
\begin{definition}[Matches Coherent In Phase Source Scenario]
  \label{def:physics:phyx-mini-0219:phyxminiproblems-problemphyxmini0219-matchescoherentinphasesourcescenario}
  \lean{PhyXMiniProblems.ProblemPhyXMini0219.MatchesCoherentInPhaseSourceScenario}
  \uses{def:physics:phyx-mini-0219:phyxminiproblems-problemphyxmini0219-speakerlabel, def:physics:phyx-mini-0219:phyxminiproblems-problemphyxmini0219-twospeakerinterferencesetup}
  Both sources are loudspeakers, mutually coherent, and emitted in phase. The common phase origin is chosen as zero cycles. This contains no microphone movement or answer-choice information
\end{definition}
\begin{proof}
  This is the declared model definition of Matches Coherent In Phase Source Scenario.
\end{proof}
\begin{definition}[Matches Problem Readouts]
  \label{def:physics:phyx-mini-0219:phyxminiproblems-problemphyxmini0219-matchesproblemreadouts}
  \lean{PhyXMiniProblems.ProblemPhyXMini0219.MatchesProblemReadouts}
  \uses{def:physics:phyx-mini-0219:phyxminiproblems-problemphyxmini0219-lengthinmeters, def:physics:phyx-mini-0219:phyxminiproblems-problemphyxmini0219-frequencyinhertz, def:physics:phyx-mini-0219:phyxminiproblems-problemphyxmini0219-twospeakerinterferencesetup}
  Numerical data explicitly stated in the problem: the speakers are 3.00 m apart, their tone is 474 Hz, and the microphone starts 3.20 m from the speaker midpoint. The initial intensity maximum is a recorded observation
\end{definition}
\begin{proof}
  This is the declared model definition of Matches Problem Readouts.
\end{proof}
\begin{definition}[Matches Supplied Figure]
  \label{def:physics:phyx-mini-0219:phyxminiproblems-problemphyxmini0219-matchessuppliedfigure}
  \lean{PhyXMiniProblems.ProblemPhyXMini0219.MatchesSuppliedFigure}
  \uses{def:physics:phyx-mini-0219:phyxminiproblems-problemphyxmini0219-signedlengthinmeters, def:physics:phyx-mini-0219:phyxminiproblems-problemphyxmini0219-twospeakerinterferencesetup}
  Coordinates read from the primary image, in meters. The speaker midpoint is the origin, the speakers are at horizontal coordinates -1.50 and 1.50, and the initial microphone lies 3.20 m vertically below the midpoint
\end{definition}
\begin{proof}
  This is the declared model definition of Matches Supplied Figure.
\end{proof}
\begin{definition}[Uses Standard Room Air Sound Speed]
  \label{def:physics:phyx-mini-0219:phyxminiproblems-problemphyxmini0219-usesstandardroomairsoundspeed}
  \lean{PhyXMiniProblems.ProblemPhyXMini0219.UsesStandardRoomAirSoundSpeed}
  \uses{def:physics:phyx-mini-0219:phyxminiproblems-problemphyxmini0219-speedinmeterspersecond, def:physics:phyx-mini-0219:phyxminiproblems-problemphyxmini0219-twospeakerinterferencesetup}
  The source does not state the air temperature or sound speed. The recorded choice 0.429 m uses the conventional room-temperature still-air calibration 343 m/s; it is therefore exposed as a separate model input rather than misattributed to the prose or figure
\end{definition}
\begin{proof}
  This is the declared model definition of Uses Standard Room Air Sound Speed.
\end{proof}
\begin{definition}[Has Physical Parameters]
  \label{def:physics:phyx-mini-0219:phyxminiproblems-problemphyxmini0219-hasphysicalparameters}
  \lean{PhyXMiniProblems.ProblemPhyXMini0219.HasPhysicalParameters}
  \uses{def:physics:phyx-mini-0219:phyxminiproblems-problemphyxmini0219-lengthquantity, def:physics:phyx-mini-0219:phyxminiproblems-problemphyxmini0219-lengthinmeters, def:physics:phyx-mini-0219:phyxminiproblems-problemphyxmini0219-frequencyinhertz, def:physics:phyx-mini-0219:phyxminiproblems-problemphyxmini0219-speedinmeterspersecond, def:physics:phyx-mini-0219:phyxminiproblems-problemphyxmini0219-pathlabel, def:physics:phyx-mini-0219:phyxminiproblems-problemphyxmini0219-twospeakerinterferencesetup}
  Positive, nondegenerate acoustic and geometric quantities
\end{definition}
\begin{proof}
  This is the declared model definition of Has Physical Parameters.
\end{proof}
\begin{definition}[Satisfies Acoustic Wave Law]
  \label{def:physics:phyx-mini-0219:phyxminiproblems-problemphyxmini0219-satisfiesacousticwavelaw}
  \lean{PhyXMiniProblems.ProblemPhyXMini0219.SatisfiesAcousticWaveLaw}
  \uses{def:physics:phyx-mini-0219:phyxminiproblems-problemphyxmini0219-lengthreadout, def:physics:phyx-mini-0219:phyxminiproblems-problemphyxmini0219-frequencyreadout, def:physics:phyx-mini-0219:phyxminiproblems-problemphyxmini0219-speedreadout, def:physics:phyx-mini-0219:phyxminiproblems-problemphyxmini0219-twospeakerinterferencesetup}
  The nondispersive acoustic-wave relation v = λ f, stated for every compatible selection of length and time units. This determines wavelength from the calibration and tone frequency but says nothing about a minimum
\end{definition}
\begin{proof}
  This is the declared model definition of Satisfies Acoustic Wave Law.
\end{proof}
\begin{definition}[Satisfies Right Move And Path Geometry]
  \label{def:physics:phyx-mini-0219:phyxminiproblems-problemphyxmini0219-satisfiesrightmoveandpathgeometry}
  \lean{PhyXMiniProblems.ProblemPhyXMini0219.SatisfiesRightMoveAndPathGeometry}
  \uses{def:physics:phyx-mini-0219:phyxminiproblems-problemphyxmini0219-lengthquantity, def:physics:phyx-mini-0219:phyxminiproblems-problemphyxmini0219-lengthreadout, def:physics:phyx-mini-0219:phyxminiproblems-problemphyxmini0219-signedlengthreadout, def:physics:phyx-mini-0219:phyxminiproblems-problemphyxmini0219-pathlabel, def:physics:phyx-mini-0219:phyxminiproblems-problemphyxmini0219-speakerforpath, def:physics:phyx-mini-0219:phyxminiproblems-problemphyxmini0219-twospeakerinterferencesetup}
  General rightward-translation and Euclidean propagation geometry. A nonnegative displacement adds to the microphone's horizontal coordinate and leaves its vertical coordinate fixed. The speaker separation, initial midpoint offset, and each labelled path are Euclidean distances. These laws hold for arbitrary displacements and contain no distinguished 0.429 m case
\end{definition}
\begin{proof}
  This is the declared model definition of Satisfies Right Move And Path Geometry.
\end{proof}
\begin{definition}[Satisfies In Phase Destructive Interference Law]
  \label{def:physics:phyx-mini-0219:phyxminiproblems-problemphyxmini0219-satisfiesinphasedestructiveinterferencelaw}
  \lean{PhyXMiniProblems.ProblemPhyXMini0219.SatisfiesInPhaseDestructiveInterferenceLaw}
  \uses{def:physics:phyx-mini-0219:phyxminiproblems-problemphyxmini0219-lengthquantity, def:physics:phyx-mini-0219:phyxminiproblems-problemphyxmini0219-lengthreadout, def:physics:phyx-mini-0219:phyxminiproblems-problemphyxmini0219-twospeakerinterferencesetup}
  For two mutually coherent, in-phase sources in the path-difference-only model, an intensity minimum occurs exactly when the magnitude of the path difference is an odd half-integral number of wavelengths. The natural order 0 is the first destructive order. This is a general governing law and does not select a displacement or an answer value
\end{definition}
\begin{proof}
  This is the declared model definition of Satisfies In Phase Destructive Interference Law.
\end{proof}
\begin{definition}[Is Positive Right Move To Minimum]
  \label{def:physics:phyx-mini-0219:phyxminiproblems-problemphyxmini0219-ispositiverightmovetominimum}
  \lean{PhyXMiniProblems.ProblemPhyXMini0219.IsPositiveRightMoveToMinimum}
  \uses{def:physics:phyx-mini-0219:phyxminiproblems-problemphyxmini0219-lengthquantity, def:physics:phyx-mini-0219:phyxminiproblems-problemphyxmini0219-lengthinmeters, def:physics:phyx-mini-0219:phyxminiproblems-problemphyxmini0219-twospeakerinterferencesetup}
  A strictly positive rightward movement reaching an intensity minimum
\end{definition}
\begin{proof}
  This is the declared model definition of Is Positive Right Move To Minimum.
\end{proof}
\begin{definition}[Is First Positive Right Move To Minimum]
  \label{def:physics:phyx-mini-0219:phyxminiproblems-problemphyxmini0219-isfirstpositiverightmovetominimum}
  \lean{PhyXMiniProblems.ProblemPhyXMini0219.IsFirstPositiveRightMoveToMinimum}
  \uses{def:physics:phyx-mini-0219:phyxminiproblems-problemphyxmini0219-lengthquantity, def:physics:phyx-mini-0219:phyxminiproblems-problemphyxmini0219-lengthinmeters, def:physics:phyx-mini-0219:phyxminiproblems-problemphyxmini0219-twospeakerinterferencesetup, def:physics:phyx-mini-0219:phyxminiproblems-problemphyxmini0219-ispositiverightmovetominimum}
  The first positive movement reaching a minimum: it reaches a minimum and no other positive minimum-producing movement has a smaller physical length. This semantic predicate contains no numerical movement or answer label
\end{definition}
\begin{proof}
  This is the declared model definition of Is First Positive Right Move To Minimum.
\end{proof}
\begin{lemma}[sound Wavelength In Meters eq]
  \label{lem:physics:phyx-mini-0219:phyxminiproblems-problemphyxmini0219-soundwavelengthinmeters-eq}
  \lean{PhyXMiniProblems.ProblemPhyXMini0219.soundWavelengthInMeters_eq}
  \uses{def:physics:phyx-mini-0219:phyxminiproblems-problemphyxmini0219-lengthinmeters, def:physics:phyx-mini-0219:phyxminiproblems-problemphyxmini0219-twospeakerinterferencesetup, def:physics:phyx-mini-0219:phyxminiproblems-problemphyxmini0219-matchesproblemreadouts, def:physics:phyx-mini-0219:phyxminiproblems-problemphyxmini0219-usesstandardroomairsoundspeed, def:physics:phyx-mini-0219:phyxminiproblems-problemphyxmini0219-satisfiesacousticwavelaw}
  The calibrated speed and the stated frequency determine the sound wavelength as 343 / 474 meters. This is an intermediate wave-law consequence, not the requested microphone movement
\end{lemma}
\begin{proof}
  The conclusion follows from the governing relations and figure readouts named in the statement of sound Wavelength In Meters eq.
\end{proof}
\begin{definition}[Answer Choice]
  \label{def:physics:phyx-mini-0219:phyxminiproblems-problemphyxmini0219-answerchoice}
  \lean{PhyXMiniProblems.ProblemPhyXMini0219.AnswerChoice}
  Labels of the four movement distances printed with the problem
\end{definition}
\begin{proof}
  This is the declared model definition of Answer Choice.
\end{proof}
\begin{definition}[displayed Move In Meters]
  \label{def:physics:phyx-mini-0219:phyxminiproblems-problemphyxmini0219-displayedmoveinmeters}
  \lean{PhyXMiniProblems.ProblemPhyXMini0219.displayedMoveInMeters}
  \uses{def:physics:phyx-mini-0219:phyxminiproblems-problemphyxmini0219-answerchoice}
  Movement distance in meters printed beside each answer label
\end{definition}
\begin{proof}
  This is the declared model definition of displayed Move In Meters.
\end{proof}
\begin{definition}[recorded Dataset Answer]
  \label{def:physics:phyx-mini-0219:phyxminiproblems-problemphyxmini0219-recordeddatasetanswer}
  \lean{PhyXMiniProblems.ProblemPhyXMini0219.recordedDatasetAnswer}
  \uses{def:physics:phyx-mini-0219:phyxminiproblems-problemphyxmini0219-answerchoice}
  Answer label recorded by the source dataset
\end{definition}
\begin{proof}
  This is the declared model definition of recorded Dataset Answer.
\end{proof}
\begin{definition}[Rounds To Displayed Move]
  \label{def:physics:phyx-mini-0219:phyxminiproblems-problemphyxmini0219-roundstodisplayedmove}
  \lean{PhyXMiniProblems.ProblemPhyXMini0219.RoundsToDisplayedMove}
  \uses{def:physics:phyx-mini-0219:phyxminiproblems-problemphyxmini0219-lengthquantity, def:physics:phyx-mini-0219:phyxminiproblems-problemphyxmini0219-lengthinmeters, def:physics:phyx-mini-0219:phyxminiproblems-problemphyxmini0219-answerchoice, def:physics:phyx-mini-0219:phyxminiproblems-problemphyxmini0219-displayedmoveinmeters}
  A derived movement rounds to a displayed three-decimal meter value when it is within half of 0.001 m of that value
\end{definition}
\begin{proof}
  This is the declared model definition of Rounds To Displayed Move.
\end{proof}
\begin{definition}[Is Unique Matching Answer]
  \label{def:physics:phyx-mini-0219:phyxminiproblems-problemphyxmini0219-isuniquematchinganswer}
  \lean{PhyXMiniProblems.ProblemPhyXMini0219.IsUniqueMatchingAnswer}
  \uses{def:physics:phyx-mini-0219:phyxminiproblems-problemphyxmini0219-lengthquantity, def:physics:phyx-mini-0219:phyxminiproblems-problemphyxmini0219-answerchoice, def:physics:phyx-mini-0219:phyxminiproblems-problemphyxmini0219-roundstodisplayedmove}
  Exactly one displayed value agrees with the derived movement
\end{definition}
\begin{proof}
  This is the declared model definition of Is Unique Matching Answer.
\end{proof}
\begin{lemma}[first Minimum Rounds To429 Meters]
  \label{lem:physics:phyx-mini-0219:phyxminiproblems-problemphyxmini0219-firstminimumroundsto429meters}
  \lean{PhyXMiniProblems.ProblemPhyXMini0219.firstMinimumRoundsTo429Meters}
  \uses{def:physics:phyx-mini-0219:phyxminiproblems-problemphyxmini0219-lengthquantity, def:physics:phyx-mini-0219:phyxminiproblems-problemphyxmini0219-twospeakerinterferencesetup, def:physics:phyx-mini-0219:phyxminiproblems-problemphyxmini0219-matchesproblemreadouts, def:physics:phyx-mini-0219:phyxminiproblems-problemphyxmini0219-matchessuppliedfigure, def:physics:phyx-mini-0219:phyxminiproblems-problemphyxmini0219-usesstandardroomairsoundspeed, def:physics:phyx-mini-0219:phyxminiproblems-problemphyxmini0219-hasphysicalparameters, def:physics:phyx-mini-0219:phyxminiproblems-problemphyxmini0219-satisfiesacousticwavelaw, def:physics:phyx-mini-0219:phyxminiproblems-problemphyxmini0219-satisfiesrightmoveandpathgeometry, def:physics:phyx-mini-0219:phyxminiproblems-problemphyxmini0219-satisfiesinphasedestructiveinterferencelaw, def:physics:phyx-mini-0219:phyxminiproblems-problemphyxmini0219-isfirstpositiverightmovetominimum, def:physics:phyx-mini-0219:phyxminiproblems-problemphyxmini0219-roundstodisplayedmove}
  Euclidean geometry and the first destructive order give a rightward movement of approximately 0.428803 m, which rounds to 0.429 m. The numerical conclusion appears only here, not in any setup, readout, or governing law
\end{lemma}
\begin{proof}
  The conclusion follows from the governing relations and figure readouts named in the statement of first Minimum Rounds To429 Meters.
\end{proof}
% --- Archon declaration topology end ---
% --- Archon physics formalization source end ---
